\documentclass[10pt,a4paper]{article}

\usepackage[margin=2cm]{geometry}
\usepackage[T1]{fontenc}
\usepackage[utf8]{inputenc}
\usepackage{lmodern}
\usepackage{amsmath,amssymb}
\usepackage[numbers,sort&compress]{natbib}
\usepackage{enumitem}
\usepackage[hidelinks]{hyperref}

\setlength{\parskip}{0.3em}
\setlength{\parindent}{0pt}
\setlist{nosep,leftmargin=1.2em}

\title{\vspace{-1.5cm}Consumer Choice in Differentiated-Product Markets:\\
The BLP Model through a Systems and Control Lens}
\author{Sara Betancur Giraldo\\\small KTH Royal Institute of Technology}
\date{\vspace{-0.6em}\today}

\begin{document}
\maketitle
\vspace{-1.2em}

\section{Scope}
Ariel Pakes' central contribution to the modelling of consumer choice is the
random-coefficients discrete-choice demand system of Berry, Levinsohn and
Pakes (BLP)~\cite{BLP1995}. BLP recovers heterogeneous preferences and firms'
markups from aggregate market data while accounting for the fact that prices
respond to demand conditions that the analyst does not observe. This note
reviews the model and the literature around it, states it mathematically, and
reads it from the perspective of systems and control theory, with the aim of
identifying where game theory can enter the research of a control group.

\section{Literature}
\textbf{Foundations.} The random-utility multinomial logit of
McFadden~\cite{McFadden1974} models a consumer choosing the alternative with the
highest utility under i.i.d.\ extreme-value noise. Its closed-form choice
probabilities come at the price of the independence of irrelevant alternatives
(IIA): substitution patterns depend only on market shares.
Berry~\cite{Berry1994} showed that with market-level data the share equations
can be \emph{inverted} to recover a mean-utility index in which the unobserved
product quality enters linearly, turning a non-linear endogeneity problem into
a linear instrumental-variables (IV) problem.

\textbf{The BLP model.} BLP~\cite{BLP1995} combined the inversion with random
coefficients on product characteristics, which breaks IIA, and with a
Bertrand--Nash pricing game on the supply side, estimating both jointly by GMM
on U.S.\ automobile data. Nevo~\cite{Nevo2001} used the framework to measure
market power in the cereal industry and became a standard reference for its
implementation, while BLP~\cite{BLP2004} added consumer-level ``micro moments''
to sharpen identification of heterogeneity.

\textbf{Identification and computation.} Berry and Haile~\cite{BerryHaile2014}
established nonparametric identification of demand from market-level data under
an index restriction and suitable instruments. On the numerical side, Dub\'e,
Fox and Su~\cite{DubeFoxSu2012} showed that loose tolerances in the inner
fixed-point loop bias estimates and proposed a constrained-optimisation (MPEC)
formulation; Conlon and Gortmaker~\cite{ConlonGortmaker2020} consolidated best
practices (instruments, integration, optimisation) in the PyBLP package.

\textbf{Dynamics and Pakes' broader programme.} Ericson and
Pakes~\cite{EricsonPakes1995} model industry evolution as a stochastic game with
entry, exit and investment, solved for a Markov-perfect equilibrium; this is the
dynamic supply-side counterpart to static BLP demand. Gowrisankaran and
Rysman~\cite{GowrisankaranRysman2012} make demand itself dynamic, with
forward-looking consumers who time durable-good purchases. Pakes, Porter, Ho
and Ishii~\cite{PPHI2015} replace point-identifying equilibrium conditions with
\emph{moment inequalities} derived from revealed preference, yielding set
estimates robust to agents' imperfect information.

\section{Mathematical Model}
Consider markets $t=1,\dots,T$, products $j=1,\dots,J_t$ and an outside good
$j=0$. Consumer $i$ obtains utility
\begin{equation}
u_{ijt} = \underbrace{x_{jt}^\top\beta - \alpha p_{jt} + \xi_{jt}}_{\delta_{jt}}
 + \underbrace{\textstyle\sum_k x_{jtk}\,\sigma_k \nu_{ik}}_{\mu_{ijt}(\theta_2)}
 + \varepsilon_{ijt}, \qquad u_{i0t}=\varepsilon_{i0t},
\end{equation}
with observed characteristics $x_{jt}$, price $p_{jt}$, unobserved quality
$\xi_{jt}$, taste shocks $\nu_i\sim F$ and $\varepsilon_{ijt}$ i.i.d.\ type-I
extreme value. Integrating over consumers gives the market shares
\begin{equation}
s_{jt}(\delta_t;\theta_2)=\int \frac{\exp(\delta_{jt}+\mu_{ijt})}
{1+\sum_{k=1}^{J_t}\exp(\delta_{kt}+\mu_{ikt})}\,dF(\nu_i).
\end{equation}
\emph{Inversion.} For given $\theta_2$, the observed shares $\mathcal{S}_t$ are
matched by the unique fixed point of
\begin{equation}
\delta_t^{h+1}=\Phi(\delta_t^h)=\delta_t^h+\ln \mathcal{S}_t-\ln s_t(\delta_t^h;\theta_2),
\end{equation}
which is a contraction~\cite{BLP1995}. \emph{Estimation.} Since prices are
correlated with $\xi$, one recovers
$\xi_{jt}(\theta)=\delta_{jt}(\theta_2)-x_{jt}^\top\beta+\alpha p_{jt}$ and
solves $\min_\theta\,\bar g(\theta)^\top W \bar g(\theta)$ with
$\bar g(\theta)=\frac1N\sum_{jt} Z_{jt}\,\xi_{jt}(\theta)$, where the
instruments $Z$ (cost shifters, rivals' characteristics) satisfy
$\mathbb{E}[\xi_{jt}\mid Z_{jt}]=0$.
\emph{Supply.} Firm $f$ owning products $\mathcal{F}_f$ maximises
$\sum_{j\in\mathcal F_f}(p_j-mc_j)\,s_j(p)$. The Bertrand--Nash first-order
conditions give
\begin{equation}
p = mc + \Delta(p)^{-1}s(p),\qquad
\Delta_{jk}=-\frac{\partial s_k}{\partial p_j}\,\mathbf{1}\{j,k\text{ same owner}\},
\end{equation}
so observed prices and estimated demand reveal marginal costs and markups.

\section{A Systems and Control Reading}
\begin{itemize}
\item \textbf{Price endogeneity is closed-loop identification.} Firms set the
input $p$ through a feedback law, (4), that reacts to the disturbance $\xi$. As
in closed-loop system identification~\cite{Ljung1999}, regressing output on
input is biased; cost shifters play the role of an external reference signal,
and BLP instruments are the econometric analogue of IV methods in system
identification.
\item \textbf{Logit choice is entropy-regularised optimisation.} The logit share
vector solves $\max_{s\in\Delta}\langle s,u\rangle+H(s)$ with $H(s)=-\sum_j
s_j\ln s_j$, i.e.\ a softmax (smoothed best response). Mixed logit is a
heterogeneous population of such agents.
\item \textbf{Inversion and equilibrium are fixed-point problems.} (3) is a
contraction iteration whose modulus determines convergence speed; the inner
tolerance issue of~\cite{DubeFoxSu2012} is an error-propagation question, and
MPEC turns the nested loop into a single constrained programme, as in MPC.
\item \textbf{Recovering costs is an inverse game.} (4) inverts Nash conditions
to recover agents' objectives from observed equilibrium play, the game-theoretic
counterpart of inverse optimal control.
\item \textbf{Shares are a population state.} With shares as a state on the
simplex, the logit dynamics $\dot s=\sigma(u(s))-s$ of population
games~\cite{Sandholm2010} give equilibrium as a rest point whose stability can
be analysed with Lyapunov and passivity tools; with a continuum of
heterogeneous, forward-looking consumers the natural object is a mean-field
game~\cite{HuangMalhameCaines2006}.
\item \textbf{Dynamic models are stochastic dynamic games.} Ericson--Pakes
equilibria are feedback (Markov) Nash equilibria in the sense of
Ba\c{s}ar--Olsder~\cite{BasarOlsder1999}, and moment
inequalities~\cite{PPHI2015} parallel set-membership estimation in robust
control.
\end{itemize}

\section{Integrating Game Theory: Research Directions}
\begin{enumerate}
\item \textbf{Feedback Nash pricing with demand dynamics:} couple a dynamic
demand state (adoption, inventories, network effects) with firms' pricing as a
differential or difference game; study existence, uniqueness and stability of
feedback equilibria~\cite{BasarOlsder1999,EricsonPakes1995}.
\item \textbf{Learning in games and adaptive control:} firms that estimate BLP
demand online while pricing face a dual-control trade-off between exploration
and exploitation; convergence of best-response or gradient learning to the
Bertrand--Nash equilibrium is an open systems question.
\item \textbf{Leader--follower design:} a regulator or platform (Stackelberg
leader) steers a population of logit consumers and competing firms through
taxes, subsidies or information, linking incentive design with population-game
stability~\cite{Sandholm2010}.
\item \textbf{Identification of games from data:} closed-loop identification
theory~\cite{Ljung1999} can inform instrument design, excitation conditions and
inverse-game estimation of agents' objectives in multi-agent systems.
\item \textbf{Large-population limits:} mean-field games~\cite{HuangMalhameCaines2006}
offer a tractable model of many strategic consumers with BLP-style heterogeneity.
\end{enumerate}
In summary, BLP is at its core an estimation problem for a static Nash
equilibrium of a population game under feedback; making the dynamics,
learning and control explicit is where control theory can contribute most.

\bibliographystyle{plainnat}
{\small\bibliography{references}}

\end{document}
